\documentclass[11pt]{article}
\usepackage{amsmath,amssymb,amsthm,hyperref,geometry}\geometry{margin=1in}
\newtheorem{theorem}{Theorem}\newtheorem{proposition}{Proposition}
\title{Independent occurrence interaction for supplied effect fibres}
\date{15 November 2026}\begin{document}\maketitle
\paragraph{Status.} Alternative independent context-input theorem, with written proofs independently reviewed in the stated scope. The pure base is $\mathbf{Set}$ with \emph{all} functions. The context/content modules below are literal native comma presentations, not reflected full-bundle modules. The content hypotheses are the original O1 relative occurrence/display hypotheses, restated here. The context input is different: supplied effect fibres compose by dependent bind, not a normalized disjoint-sum rig. Standard commutative-monad traversal is an ingredient, not a novelty claim or substitute for the native proof.
\paragraph{Size.} Fix nested universes for carriers, their categories and retained-native-arrow relative pseudocoends.

\section{Independent context input}
Let $T$ be a commutative monad on sets with its finite Fubini products. Affinity is allowed but not needed in the theorem below: no arbitrary discarding equation is invoked. Let $\mathcal J$ contain identities and isomorphisms, and be closed under pullback and composition. An admitted context over $A$ is a family
\[
 P=(X_a,p_a\in TX_a)_{a\in A},\qquad
 e_AP=\coprod_a X_a,\quad\pi_P:e_AP\to A\ \in\mathcal J.
\]
Its native arrow $u:P\to P'$ is a family $u_a:X_a\to X'_a$ with $Tu_a(p_a)=p'_a$. These arrows need not be invertible or surjective. Reindex along any $f:A\to A'$ by $f^*P'=(X'_{f(a)},p'_{f(a)})$. Its extension square is the actual pullback of $\pi_{P'}$.

Unit context has singleton fibres and effect $\eta(*)$. For $R=(Y_{a,x},q_{a,x})$ over $e_AP$, define
\[
 (P\star R)_a=\coprod_{x\in X_a}Y_{a,x},\quad
 (p\star q)_a=p_a\gg=(x\mapsto q_{a,x}\gg=(y\mapsto\eta(x,y))).\tag{K}
\]
Its projection is the composite of the two admitted carrier projections, so it remains in $\mathcal J$. Native maps on composable pairs send $(x,y)$ to $(u_a x,v_{a,x}y)$; bind naturality proves effect preservation. Tuple reassociation and the monad laws prove the unit, associativity and reindexing laws on raw witnesses and native maps. Hence
\[
 M(A,B)=(e_A\downarrow B)
\]
is an independently valid context relative pseudopromonad. Its objects have arbitrary pure body $g:e_AP\to B$; composition reindexes the next context along $g$ before (K). No normalization or conditioning is used. In a nonaffine instance the unique raw map to a singleton is not automatically a native context arrow: it must still preserve the supplied effect. Context units use $\eta(*)$ and monad unit laws, not an implicit discard law.

\section{Independent content input}
Let $S(B)=\coprod_{c\in C} B^{I_c}$ be a finitary cartesian polynomial monad, with $S(\mathcal J)\subseteq\mathcal J$. In particular its unit has one position and its multiplication has specified addressed flattening bijections
\[
 I_{m(c,(c_i))}\cong\coprod_{i\in I_c} I_{c_i},
\]
coherent by the polynomial monad laws. Supply an already valid contentad $p:\mathcal L\to\mathbf{Set}$, extension $d$, unit $U$, and result-dependent accumulation $Q\diamond R$ with $R\in\mathcal L_{dQ}$. Thus $N(A,B)=(A\downarrow d_B)$ is independently a relative pseudopromonad. Require:
\begin{enumerate}
\item Given global opcartesian transport, and genuine cartesian lifts over every $f\in\mathcal J$, universal against arbitrary base arrows.
\item A natural occurrence map $\nu:d\Rightarrow Sp$, on all total-category arrows, commuting with content unit and accumulation via the unit and multiplication of $S$.
\item For every admitted cartesian lift $c_{f,Q}:f^*Q\to Q$, the canonical display map
\[
 d(f^*Q)\longrightarrow dQ\times_{SB}SX,
 \quad z\mapsto(d c_{f,Q}(z),\nu_{f^*Q}(z))
\]
is a bijection.
\item Content unit and accumulation preserve these cartesian arrows, including the composable-pair lifting structure for accumulation.
\item In every pullback square with vertical maps in $\mathcal J$ and arbitrary horizontal maps $u,v$, the canonical mate $v_!f^*\to f'^*u_!$ is invertible.
\end{enumerate}
These are supplied one-sided content/display conditions, not assumptions of a mixed crossing. No claim that an arbitrary contentad has them is made.

\section{Construction on the full native carrier}
For $P$ over $B$ and $Q\in\mathcal L_B$, write $\nu_Q(z)=(c,(b_i)_{i\in I_c})$. Define a context over $dQ$ by
\[
 \widehat P_z=\prod_{i\in I_c} X_{b_i},\qquad
 \widehat p_z=\bigotimes_{i\in I_c}p_{b_i},\qquad
 \widehat Q=\pi_P^*Q.\tag{X}
\]
The raw carrier of $\widehat P$ is canonically $dQ\times_{SB}S(e_BP)$. Its projection is a pullback of $S\pi_P$, so it lies in $\mathcal J$. The inverse content-display map gives the crossing body
\[
 e_{dQ}\widehat P\ \xrightarrow{\cong}\ d(\pi_P^*Q).\tag{B}
\]
It retains both the original $z$ and its entire addressed decorated occurrence. Nullary occurrences have singleton context fibre and effect $\eta(*)$. Equal $b_i$ at different addresses are sampled as distinct Fubini factors; they are not identified.

For a general writer/content-then-context representative $(Q,f:A\to dQ)$ followed by $(P,g:e_BP\to D)$, pull (X) along $f$ and push $\widehat Q$ along $g$. Compose (B) with the given content opcartesian extension map. This constructs a representative of $M\odot N$ from one of $N\odot M$ on the original pure base.

\section{Native arrows and arbitrary pure middle balance}
For a context map $u:P\to P'$, its total carrier $e u:e_BP\to e_BP'$ lies over $B$. The crossed context map at $z$ is the product of the actual $u_{b_i}$; its effect equation is finite Fubini naturality. The crossed content arrow
\[
 (e u)_!(\pi_P^*Q)\longrightarrow\pi_{P'}^*Q
\]
is obtained by factoring through the target cartesian lift and then the source opcartesian lift. It need not be invertible. The display coordinates prove its body square: the original content point agrees, and its decorated $S(e_BP')$ label is the coordinatewise image.

For a content arrow $a:Q\to Q'$ over $B$, naturality says $\nu_{Q'}d a=\nu_Q$. Corresponding occurrences therefore have exactly the same shape, addresses and labels. The crossed context map is identity on the family in (X), and the content map is $\pi_P^*a$. For arrows over a nonidentity base map, labels change by that base map while addresses remain; input reindexing supplies the matching fibres. Arbitrary address duplication is not part of this hypothesis and no relevance/copyability law is required.

Identity and composition of native content maps are preserved by cartesian functoriality. For native context maps both universal factorizations compose, while product maps compose coordinatewise; the same argument proves mixed naturality. Pure-middle balance along $u:B\to B'$ uses the actual pullback of the context carrier projection and the supplied invertible Beck--Chevalley mate. On contexts it is the canonical reindexing of the same addressed product and effect; on contents it is that mate; the displayed body commutes since its two display coordinates agree. Mate-pasting and pullback identities supply the balance coherence. This establishes a native module transformation
\[
 \tau:N\odot M\longrightarrow M\odot N,
\]
not merely an object-level law on final effects.

\section{The four pseudodistributive comparisons}
For the context unit, each occurrence carries a singleton witness with unit effect, so (X) is the unit context; the cartesian lift along an identity gives the content comparison. For the content unit, its unit body has one occurrence; the crossed context is the original context. Preservation of cartesian arrows by $U$ supplies the corresponding content-unit comparison.

For context multiplication, at each $z$ the essential raw bijection is
\[
 \prod_{i\in I_c}\Bigl(\coprod_{x_i\in X_{b_i}}Y_{b_i,x_i}\Bigr)
 \ \cong\
 \coprod_{(x_i)\in\prod_iX_{b_i}}\prod_{i\in I_c}Y_{b_i,x_i}.\tag{KX}
\]
Its effect equation is exactly
\[
 \bigotimes_i(p_i\gg=q_i)
 =\Bigl(\bigotimes_i p_i\Bigr)\gg=
       \bigl((x_i)\mapsto\bigotimes_i q_i(x_i)\bigr),
\]
where each continuation retains its selected $x_i$. This is commutative-monad Fubini/bind interchange. The content comparison is the canonical composition-of-cartesian-lifts isomorphism for the two carrier projections, followed by the given transports. Thus no nonexistent normalized coproduct operation is needed.

For content multiplication, the occurrence axiom gives the flattening bijection from nested addressed positions to leaf positions. Both crossed contexts have the product of the same effects over these positions, identified by finite Fubini associativity and symmetry. Preservation of cartesian arrows by content accumulation gives the content-grade comparison. The two display coordinates give the body equation. These four comparisons are invertible, although the crossing functor and general native cells need not be.

\begin{theorem}[Alternative effect-fibre interaction]
Under the independent hypotheses above, the crossing and four comparisons form a full pseudodistributive law. Consequently the literal comma family
\[
 H(A,B)=(e_A\downarrow d_B)\simeq M\odot N
\]
has the induced relative pseudopromonad multiplication retaining native context/content arrows and arbitrary pure endpoint squares.
\end{theorem}
\begin{proof}[Coherence and normalization]
Every raw context comparison is an addressed tuple map built from (KX), polynomial flattening, units, or reindexing. Any two routes in a distributive coherence diagram give the same raw address map: flattening coherence is the cartesian polynomial monad law, and dependent-tuple regrouping is substitution associativity. Its effect equation follows from commutative-monad coherence and bind associativity. All native context arrows are actual maps on these full raw tuples.

The content comparisons are the canonical cartesian/opcartesian factorizations, unit/accumulation preservation comparisons and Beck--Chevalley mates. Their diagrams commute by uniqueness of universal factorizations and mate-pasting, with the independently supplied contentad coherence. To compare body arrows, compose with the appropriate display bijection: the original-content coordinate agrees by those universal diagrams; the decorated occurrence coordinate agrees by the raw address calculation. Joint injectivity of the display then proves equality of the bodies. This establishes equality of full native cells, not just their effects.

Finally retained comma normalization uses $[B;(P,g),(Q,h)]\mapsto(P,Q,hg)$ and section $[e_AP;(P,\mathrm{id}),(Q,f)]$. On a native $(\alpha,\beta)$, apply the native content arrow, balance along $e\alpha$, then the native context arrow, exactly retaining both noninvertible components. Pure-middle balance gives the inverse comparison and its triangles. There is no reflective endpoint normalization here because the pure base is Set itself. Transporting the standard composite multiplication through this equivalence proves the stated conclusion.
\end{proof}

\section{Scope and comparison}
This theorem substitutes a dependent-bind contextad for the original O1 rig-valued augmented-context input, while retaining its stated content/display/BC input. It does not derive those content assumptions or its marked necessity/converse and parameter results. It is not obtained by pretending normalized measures have a disjoint-sum rig. The finite occurrence effect traversal is established algebra; the programme-level construction here is the actual category-valued independent context-input comparison with full native maps and display coherence, without a novelty claim.

Neither infinite occurrence products nor a recursive-path interpretation follow. The previously proved guarded-prefix construction is a separate unary result; extending this multi-occurrence law through guarded recursion requires a new compatible-prefix/occurrence theorem, not an appeal to finitary traversal alone. Shared-world correlated effects require explicitly retained shared data and are not made independent by (X).
\section*{Source comparison}
The one-sided content assumptions, retained comma normalization and universal display-coherence strategy are those of the programme's \emph{General Independent Interaction Criterion}, especially its relative content input, generic crossing, native-cell and pure-middle-balance sections. The present change is the context input (K) and its effect-preserving native maps; that difference must be checked rather than inferred from the old rig construction. The old source already cites Fujii, Tsai, Montacute and Hasuo, \emph{Monads and Distributive Laws in Substructural Contexts}, for the established structural-rule perspective: \url{https://doi.org/10.4230/LIPIcs.LICS.2026.45}. Kock, \emph{Commutative monads as a theory of distributions}, Theorem9.2, supplies the standard Fubini/bind algebra: \url{https://tildeweb.au.dk/au76680/CMTD.pdf}. These are reused foundations, not newly discovered gaps in the literature. A bounded literature check does not establish novelty of the alternative context-input theorem.
\end{document}
